% authority: science-draft
% status: draft/control
% route: ontology-law-research-packet
% trigger_classification: derivation_critical_missing_source_law
% adoption_status: blocked_adoption_open_continuation
\documentclass[11pt]{article}
\usepackage[margin=1in]{geometry}
\usepackage{amsmath,amssymb,amsthm}
\usepackage{booktabs,longtable,array}
\usepackage[T1]{fontenc}
\usepackage{lmodern}
\usepackage{enumitem}
\usepackage[hidelinks]{hyperref}

\newtheorem{definition}{Definition}
\newtheorem{theorem}{Theorem}
\newtheorem{proposition}{Proposition}
\newtheorem{corollary}{Corollary}
\newtheorem{remark}{Remark}
\newcommand{\Core}{\operatorname{Core}}
\newcommand{\Aut}{\operatorname{Aut}}
\newcommand{\Nat}{\operatorname{Nat}}
\newcommand{\Fix}{\operatorname{Fix}}
\newcommand{\Set}{\mathsf{Set}}
\newcommand{\EqSrc}{\mathsf{EqSrc}}

\title{V21 P2-T02 EqSrc Natural-Selector\\
Fixed-Point Condition v1}
\author{Ontology Formalizer control packet RT-20260720-015}
\date{2026-07-20}

\begin{document}
\maketitle

\section{Control status and exact burden}

This artifact executes only v21 work item \texttt{P2-T02}. Its identifier is
\texttt{EQSRC-NATURAL-SELECTOR-}\allowbreak
\texttt{FIXED-POINT-CONDITION-V1}. The exact burden is
to define deterministic source-natural selection on the already declared
P2-T01 core and derive the necessary automorphism fixed-point condition. The
P2-T01 objects, arrows, admission receipts, certified isomorphisms,
eligibility predicate, choice functor, relation interface, and automorphism
actions remain unchanged.

The record is \texttt{draft/control}, \texttt{proposal-only}, and
\texttt{source-extension data}. Current ontology does not derive the P2-T01
source category, its eligible-choice functor, a natural selector, or a
physical gauge groupoid. Adoption is
\texttt{blocked\_adoption\_open\_continuation}. The theorem below is exact
relative to the declared proposal domain; it is not ontology adoption,
physical covariance, empirical recovery, general \(\EqSrc\), or benchmark
progress.

\section{Frozen input domain}

Let
\[
  \mathcal G_\Sigma=\Core(\mathcal C_\Sigma)
\]
be the certified-isomorphism core defined by P2-T01. Let
\[
  S_\chi:\mathcal G_\Sigma\longrightarrow\Set
\]
be its eligible-choice functor. For each certified isomorphism
\(u:X\to Y\), the map \(S_\chi(u)\) is exactly the action already declared
by P2-T01. This packet does not remove an arrow, add an arrow, modify
\(\chi\), alter \(S_\chi\), or redefine an action after inspecting a desired
choice or induced relation.

Let \(1_{\mathcal G_\Sigma}:\mathcal G_\Sigma\to\Set\) be the constant
singleton functor. It sends every object to \(\{\star\}\) and every arrow to
the identity of that singleton.

\begin{definition}[Deterministic source-natural selector]
A deterministic source-natural selector is a natural transformation
\[
  \sigma:1_{\mathcal G_\Sigma}\Longrightarrow S_\chi.
\]
Equivalently, it is a family \(\sigma_X\in S_\chi(X)\) satisfying
\[
  S_\chi(u)(\sigma_X)=\sigma_Y
  \qquad\text{for every certified isomorphism }u:X\to Y.
\]
The selector space is
\[
  \operatorname{Sel}_\chi
  :=\Nat(1_{\mathcal G_\Sigma},S_\chi).
\]
It is equivalently the space of natural sections of the unique natural map
\(S_\chi\to1_{\mathcal G_\Sigma}\). The word ``deterministic'' means only
single-valued; it asserts neither computability nor a physical determinism
principle.
\end{definition}

\begin{definition}[Invariant fixed locus and relation image]
For each object \(X\), define
\[
 F_X:=S_\chi(X)^{\Aut_{\mathcal G_\Sigma}(X)}
 =\{d\in S_\chi(X):S_\chi(g)d=d\text{ for every }
 g\in\Aut_{\mathcal G_\Sigma}(X)\}.
\]
If the P2-T01 natural relation map \(K_X\) is present, define separately
\[
  \mathcal R_X:=K_X(F_X).
\]
The sets \(F_X\) and \(\mathcal R_X\) answer different questions: the first
concerns selected source choices, while the second concerns their induced
relation images.
\end{definition}

\section{Necessary fixed-point theorem}

\begin{theorem}[Necessary automorphism fixed-point condition]
If \(\sigma\in\operatorname{Sel}_\chi\), then
\[
  \sigma_X\in F_X
\]
for every object \(X\) of \(\mathcal G_\Sigma\).
\end{theorem}

\begin{proof}
Fix an object \(X\) and an arbitrary
\(g\in\Aut_{\mathcal G_\Sigma}(X)\). Naturality of \(\sigma\) gives
\[
  S_\chi(g)\circ\sigma_X
  =\sigma_X\circ1_{\mathcal G_\Sigma}(g).
\]
The singleton functor sends \(g\) to the identity. Evaluating the displayed
equality at \(\star\) therefore yields
\[
  S_\chi(g)(\sigma_X)=\sigma_X.
\]
Because \(g\) was arbitrary, \(\sigma_X\) lies in \(F_X\).
\end{proof}

\begin{corollary}[Empty-fixed obstruction]
If \(F_X=\varnothing\) for any object \(X\), then no natural selector exists
on the connected component containing \(X\), and consequently
\(\operatorname{Sel}_\chi=\varnothing\) on the full groupoid.
\end{corollary}

\begin{proof}
A selector on the component, or on the full groupoid, would restrict at
\(X\) to an element of \(F_X\), contradicting emptiness.
\end{proof}

\begin{remark}[Singleton and multiple fixed loci]
If \(|F_X|=1\), every selector, if one exists, is forced at \(X\). This is a
local uniqueness statement, not an existence proof. If \(|F_X|>1\),
automorphism naturality alone does not distinguish a unique value at \(X\).
It does not follow that the full groupoid has multiple selectors: another
component may have an empty fixed locus, or additional coherence may fail.
Likewise \(|F_X|>1\) can coexist with \(|\mathcal R_X|=1\) when \(K_X\) is
constant or noninjective on \(F_X\).
\end{remark}

On a connected groupoid component, the prior P1-T04 artifact proves that
evaluation at a representative gives a bijection from natural sections to
the representative fixed locus. This packet cites that result and does not
relabel it as a new complete general theorem. P2-T03 retains the general
no-selector and nonuniqueness theorem, global component quantifiers, any
choice-set assumptions for infinitely many components, and the complete
relation-image classification. Extending from the core to noninvertible
arrows also requires additional coherence equations not proved here.

\section{Added source data and stabilizers}

Let \(\Sigma^+\) be an enriched proposal signature and
\[
 U:\mathcal G_{\Sigma^+}\longrightarrow\mathcal G_\Sigma
\]
a target-neutral reduct functor. Let \(X^+\) reduce to \(X\). The induced
map on automorphisms is
\[
 U_*:\Aut_{\Sigma^+}(X^+)\longrightarrow\Aut_\Sigma(X),
 \qquad H_{X^+}:=\operatorname{im}U_*.
\]

\begin{proposition}[Added-mark stabilizer and fixed-locus inclusion]
Suppose \(X^+\) contains an added mark \(m\) preserved by every enriched
automorphism. Then
\[
 H_{X^+}\leq\operatorname{Stab}_{\Aut_\Sigma(X)}(m).
\]
If, additionally, the eligible-choice set is unchanged and the enriched
action is exactly the restriction of the base action to \(H_{X^+}\), then
\[
 S_\chi(X)^{\Aut_\Sigma(X)}
 \subseteq S_\chi(X)^{H_{X^+}}.
\]
\end{proposition}

\begin{proof}
Functoriality of \(U\) makes \(U_*\) a group homomorphism. Every enriched
automorphism preserves \(m\); hence every reduct in its image lies in the
stabilizer of \(m\). For the second statement, a point fixed by every element
of the base group is fixed by every element of its subgroup \(H_{X^+}\).
\end{proof}

Equality with the full mark stabilizer requires every mark-preserving reduct
automorphism to lift; identifying the enriched automorphism group with that
stabilizer also requires faithfulness on automorphisms. The inclusion can be
strict, and reducing a group need not create a fixed point or a unique one.
If the extension changes the choice space, no fixed-locus comparison follows
without an explicit equivariant comparison map.

A pure presentation or serialization relabeling adds no scientific structure.
An eliminable abbreviation is a
\texttt{conservative\_definitional\_extension} only when an explicit
elimination proof exists. A non-eliminable root, line, orientation, order,
grading, action, name, random seed, measure, normalization, response token,
or selector-relevant mark is a
\texttt{new\_ontology\_primitive\_candidate} with status
\texttt{draft/control}, \texttt{proposal-only}, and
\texttt{source-extension data}. If the mark is the desired choice itself,
the construction is circular. None of these cases is, by mathematics alone,
spontaneous, dynamical, or physical symmetry breaking.

\section{Finite controls}

The task-local validator enumerates the following controls.

\begin{longtable}{>{\raggedright\arraybackslash}p{0.22\textwidth}
                  >{\raggedright\arraybackslash}p{0.30\textwidth}
                  >{\raggedright\arraybackslash}p{0.38\textwidth}}
\toprule
Control & Exact count & Bounded consequence \\
\midrule
\endhead
\(\mathrm{GL}(2,2)\) lines & Six automorphisms act on three lines; no
common fixed line. & Empty fixed locus obstructs a selector on this one-object
component. \\
Directed \(C_4\) roots & Four rotations act on four roots; no common fixed
root. & Empty fixed locus obstructs a selector on this one-object component. \\
Ordered four-chain & The source-declared eligible set is the singleton global
minimum and the structural automorphism group is trivial. & Exactly one fixed
choice and one selector exist in this proposal control. \\
Trivial two-choice object & Both choices are fixed; injective \(K\) has two
images and constant \(K\) has one. & Naturality alone does not select uniquely;
relation-level uniqueness is independent. \\
Marked \(\mathrm{GL}(2,2)\) line & The image stabilizer has order two and
fixes exactly the marked line. & Added data can enlarge a fixed locus, but the
mark remains loaded source structure. \\
Marked \(C_4\) root & The stabilizer is trivial and therefore fixes all four
unchanged eligible roots. & Symmetry reduction alone does not imply uniqueness. \\
Noninvertible coherence guard & Two objectwise fixed assignments exist but
only one respects the declared noninvertible arrow. & Automorphism fixedness
does not establish full-category coherence. \\
\bottomrule
\end{longtable}

A disconnected counterexample guards one further overread: the disjoint union
of a trivial one-object groupoid with two fixed choices and a one-object
\(C_2\)-groupoid acting freely on two choices has no global selector. Thus
multiple fixed points on one component do not imply multiple selectors on the
full groupoid.

The controls are mathematical enumerations. They are not observations,
predictions, physical gauge evidence, or exact-GR recovery evidence.

\section{Claim boundary and next ownership}

The proved payload is only the necessary fixed-point theorem and the bounded
added-data lemma on the unchanged P2-T01 proposal domain. Structural
automorphisms are defined relative to \(\Sigma\); they are not physical gauge
transformations without separately derived or adopted states, dynamics or
constraints, observables, and a redundancy interpretation.

No object, arrow, eligibility rule, \(K\)-map, or automorphism action was
selected after observing a desired result. No exact candidate family was
reopened or rejected. No source law or canonical ontology change is adopted.
General \(\EqSrc\), \(M_{\rm src}\), \(g_{\rm eff}\), matter coupling,
Einstein equations, benchmark promotion, Gate Chair authority, publication,
and completed derivation remain blocked. Both scientific ledgers remain
unchanged.

P2-T03 owns the complete no-selector and nonuniqueness theorem. P2-T06 owns an
independent audit for hidden choice, target import, post-hoc category
engineering, and false conservativity. P2-T07 owns independent stress under
weakened or strengthened categories, altered actions, noninvertible arrows,
and theorem-boundary countermodels. Protected human gates retain ontology or
source-law adoption, physical-gauge identification, benchmark promotion, and
Gate Chair decisions.

\section{Source record}

The formal terminology follows the registered P1-T05 primary-literature
review and its source-specific transfer limits. Natural transformations are
used in the standard categorical sense of Eilenberg and Mac Lane (1945).
Project-specific theorem and boundary authority comes from the registered
P1-T04 target, the P2-T01 source-category definition, the v21 plan, and
handoff-0784, not from analogy alone.

\begin{thebibliography}{9}
\bibitem{eilenberg-maclane-1945}
Eilenberg, S., \& Mac Lane, S. (1945). General theory of natural
equivalences. \emph{Transactions of the American Mathematical Society, 58},
231--294. \url{https://doi.org/10.1090/S0002-9947-1945-0013131-6}
\end{thebibliography}

\end{document}
